\documentclass[11pt]{article}

\usepackage[margin=1in]{geometry}
\usepackage[utf8]{inputenc}
\usepackage{amsmath,amssymb,amsthm,mathtools}
\usepackage{microtype}
\usepackage{booktabs,tabularx,array}
\usepackage{enumitem}
\usepackage[round,authoryear]{natbib}
\usepackage{titlesec}
\usepackage[colorlinks,allcolors=black]{hyperref}

% ---------- layout ----------
\titleformat{\section}{\large\bfseries}{\thesection}{0.6em}{}
\titlespacing*{\section}{0pt}{1.6ex plus .4ex minus .2ex}{0.7ex}
\titleformat{\paragraph}[runin]{\normalsize\bfseries}{}{0pt}{}[.]
\titlespacing*{\paragraph}{0pt}{1.1ex plus .3ex minus .2ex}{0.6em}
\setlist{nosep,leftmargin=2.1em}
\setlength{\bibsep}{2pt}
\newcolumntype{Y}{>{\raggedright\arraybackslash}X}
\setlength{\intextsep}{8pt plus 2pt minus 2pt}
\setlength{\abovecaptionskip}{4pt}
\setlength{\belowcaptionskip}{0pt}

% ---------- theorem environments ----------
\newtheoremstyle{tight}{5pt}{5pt}{\itshape}{}{\bfseries}{.}{0.5em}{}
\theoremstyle{tight}
\newtheorem{proposition}{Proposition}
\newtheorem{assumption}{Assumption}
\renewcommand{\theassumption}{A\arabic{assumption}}

% ---------- macros ----------
\newcommand{\E}{\mathbb{E}}
\newcommand{\R}{\mathbb{R}}
\newcommand{\Sset}{\mathcal{S}}
\newcommand{\one}{\mathbf{1}}
\DeclareMathOperator{\med}{med}

\begin{document}
\begin{center}
{\Large\bfseries Learning a Common Index from Outcomes Measured\\[0.15em] on Non-Comparable Scales}\\[0.5em]
Project proposal \quad$\cdot$\quad Nicolas Abbate \quad$\cdot$\quad October 5, 2026
\end{center}
\vspace{-0.6em}

% =====================================================================
\section{Motivation}
% =====================================================================

Aerial imagery now covers most of the United States at building-level resolution and annual
frequency, while the American Community Survey reports five-year averages for census tracts
with margins of error that are wide for small areas. A model that reads wealth from imagery
could close that gap. It would help target programs at a geography the survey cannot support,
sharpen small-area estimates by combining the two signals, and measure change between survey
rounds. It would also make impact evaluation possible where no survey covers the affected
places or years: after a rezoning or a transit extension, the units of interest are blocks and
the window is shorter than a survey cycle.

Remote wealth estimation has a fundamental problem that remains unresolved. Detection
models scale because a building is a building in Dallas and in Lagos, so every labeled image
is one more example of a single relationship. Wealth labels lack that property. Surveys report
wealth in dollars, and the same visual bundle carries different dollar values in two cities or
in two years, because each local wealth distribution is shifted, stretched and reshaped by
forces unrelated to the built environment. Pooling dollar labels across places and periods
therefore mixes incompatible relationships, and a new city can add an inconsistent example
rather than another instance of the same one.\footnote{As an illustration, \citet{khachiyan2022} report $R^2\approx 0.90$
for income levels pooled over 2000 and 2010 and $R^2\approx 0.40$ for the 2000--2010 change. While they do not work with wealth, 
the same label distortion problem arises in their setting.}

One reading of this evidence is that the map from features to wealth is inherently local, so
every city and year needs its own model. We read it differently: the locality is in the dollar
labels, not in the ordering. Wealth buys amenities that look much the same everywhere---space,
light, quiet, greenery, quality of construction, height, centrality. Because of this, we believe 
that there exists a global relation between the features and wealth; what is not global, and changes
from city to city and year to year, is the number of dollars attached to each feature bundle.

The problem is not specific to imagery or to wealth. It arises wherever outcomes are comparable
within groups but not across them: a fixed income affords different housing across cities, exams
vary in difficulty, and credit ratings depend on the rating system. The same structure appears
in preference and reward data with annotator-specific scales, test scores across forms and
graders, credit risk across vintages, severity graded from medical images across sites and
scanners, and hedonic quality across markets whose price levels differ. In each case covariates
determine a shared unobservable index that every group reports through its own unknown monotone rescaling.

This project studies how to learn this latent shared index---or equivalently, learning the function 
that maps the covariates to the index---, without estimating the group-specific rescalings, using external data
on . Within a group, 
the labels order units as the index does. Across
groups,  known from outside the labels to be nearly the same, such as one location
photographed in two years with no structural change, must receive nearly the same index, a
restriction that uses no label. We combine them in one training criterion: a pairwise ranking
loss within groups plus a scale-free penalty on how much the score moves over such pairs. We
then prove consistency of the estimator under three assumptions on the comparison errors, and
run a Monte Carlo study of when each version works, when it fails, and whether the unlabeled
pairs improve prediction.

% =====================================================================
\section{Framework}
% =====================================================================

Units are indexed by $i$, groups by $c$ (cities) and dates by $t$; the pair $(c,t)$ is a
\emph{cell}. We observe a label $W_{ict}\in\R$ and covariates $X_{ict}$, an image tile in the
application. The model is
\begin{equation}\label{eq:model}
  W_{ict} \;=\; \xi_{ct}(X_{ict}) + e_{ict},
  \qquad
  \xi_{ct} \;=\; \psi_{ct}\circ f .
\end{equation}
The \emph{global function} $f$ maps covariates into a latent index and is the same in every
cell. The \emph{cell-specific function} $\psi_{ct}$ converts the index into the units in which
the cell reports it (dollars in city $c$ in year $t$) and carries that cell's level, spread and
shape. A model fitted cell by cell learns $\xi_{ct}$, which says nothing about a cell without
labels. The goal is to estimate $f$ without estimating $\psi_{ct}$.

Let $\Sset$ be a set of \emph{stable pairs} $(x,x')$: covariate values observed in two
different cells and selected, without reference to the labels, as nearly the same unit.

\begin{assumption}[Monotonicity]\label{a:mono}
$\psi_{ct}$ is strictly increasing for every $(c,t)$.
\end{assumption}
\begin{assumption}[Invariance]\label{a:inv}
$f$ does not depend on $c$ or $t$.
\end{assumption}
\begin{assumption}[Regularity]\label{a:lip}
There is a factorization $f=g\circ\phi$, with the coordinates of $H=\phi(X)\in\R^K$
standardized, such that $|g(h)-g(h')|\le L\,\|h-h'\|$ on the support of $H$, and every stable
pair $(x,x')\in\Sset$ satisfies $\|\phi(x')-\phi(x)\|\le\delta$. The index is not constant
within cells.
\end{assumption}

\ref{a:inv} is the substantive assumption: all variation across cells is assigned to
$\psi_{ct}$. In \ref{a:lip}, $H$ is the embedding before the final scalar layer of a network,
and $L$ and $\delta$ matter only through their product.

\begin{proposition}[Pooled levels estimate a mixture]\label{p:mix}
Suppose \eqref{eq:model} holds with $\E[e_{ict}\mid X_{ict},c,t]=0$ and $\E|W|<\infty$, and $W$ is
regressed on $X$ pooling all cells. Then
\[
  \E[W\mid X=x] \;=\; \sum_{c,t}\Pr(c,t\mid X=x)\;\psi_{ct}\bigl(f(x)\bigr).
\]
\end{proposition}

The weights reflect which cells are in the sample, and they can reverse the ranking. If public
housing towers appear only in New York and single-family houses mostly in Dallas, a large enough
dollar gap between the cities makes the pooled fit rank the towers above the houses, although
they are poorer in each city: the model has learned that public housing signals New York.

\begin{proposition}[What the data reveal about $f$]\label{p:reveal}
Let \ref{a:mono}--\ref{a:lip} hold.
\begin{enumerate}[label=\textup{(\roman*)}]
\item \textup{Ordering within a cell.} For every cell $(c,t)$ and any $x,x'$:\;
      $\xi_{ct}(x')>\xi_{ct}(x)\iff f(x')>f(x)$.
\item \textup{Stability across cells.} For every stable pair $(x,x')\in\Sset$, with $x$ observed
      in cell $(c,t)$ and $x'$ in cell $(c',t')\neq(c,t)$:\; $|f(x')-f(x)|\le L\delta$.
\end{enumerate}
\end{proposition}
\begin{proof}
(i) $\xi_{ct}=\psi_{ct}\circ f$ and $\psi_{ct}$ is strictly increasing by \ref{a:mono}.
(ii) By \ref{a:inv} the same $f=g\circ\phi$ applies in both cells, so
$|f(x')-f(x)|\le L\,\|\phi(x')-\phi(x)\|\le L\delta$ by \ref{a:lip}.
\end{proof}

Part (ii) holds for any two cells, so a stable pair may join two dates of one group
($c=c'$, $t\neq t'$) or two different groups ($c\neq c'$). In satellite and aerial imagery only
the first kind arises on its own: one location, two years, no recorded structural change.
Nothing physical ties a tile in one city to a tile in another, so cross-city pairs require an
already pretrained model, in whose representation near neighbours can be defined before training.

Proposition~\ref{p:reveal} establishes what is observable about $f$: orderings within cells and
near-equalities across cells, neither involving $\psi_{ct}$. But the labels are $W$, not
$\xi_{ct}(X)$: whether a noisy comparison still reveals the ordering of $f$ depends on the
errors, and the assumption made about them determines the loss. For two units $i\neq j$
of one cell write $y_{ij}=\one\{W_j>W_i\}$, $p_{ij}=\Pr(y_{ij}=1\mid X_i,X_j,c,t)$ and
$\Delta_{ij}s=s(X_j)-s(X_i)$ for a score function $s$.

\begin{assumption}[Comparison errors, one of three versions]\label{a:err}\leavevmode
\begin{enumerate}[label=\textup{(\alph*)}]
\item \textup{Median zero.} Given $(X_i,X_j,c,t)$, $e_i-e_j$ has median zero and a distribution
      function strictly increasing at zero. By Proposition~\ref{p:reveal}\textup{(i)},
      $p_{ij}>\tfrac12\iff\Delta_{ij}f>0$.
\item \textup{Heterogeneous logistic.} $p_{ij}=\Lambda(\rho_{ct}\,\Delta_{ij}f)$ with
      $\rho_{ct}>0$ and $\Lambda(z)=1/(1+e^{-z})$.
\item \textup{Homogeneous logistic.} Version \textup{(b)} with $\rho_{ct}=\rho$ in every cell.
\end{enumerate}
\end{assumption}

Version (a) restricts only the sign of $p_{ij}-\tfrac12$ and lets the noise depend on the
features; (b) and (c) also fix its shape, with and without a reliability that differs by
cell.\footnote{Versions (b) and (c) place the noise on the scale of the index. They hold when
$W=\Psi_{ct}(f(X)+\varepsilon)$ with logistic differences of $\varepsilon$, a model whose
regression form is \eqref{eq:model} if $\varepsilon$ is independent of $X$ within a cell; with
logistic differences of the additive $e$ they hold when $\psi_{ct}$ is affine.} All three lead to one criterion,
\begin{equation}\label{eq:crit}
  Q_k(s,\eta) \;=\; Q^{\mathrm{cross}}_k(s,\eta)
  \;+\;\lambda_t\,\frac{S_t(s)}{V(s)}
  \;+\;\lambda_c\,\frac{S_c(s)}{V(s)},
  \qquad k\in\{\mathrm{a},\mathrm{b},\mathrm{c}\},
\end{equation}
where $S_t(s)$ and $S_c(s)$ are the mean squared score differences $\E[(s(x')-s(x))^2]$ over
temporal and cross-group stable pairs, and $V(s)=\tfrac12\E[(\Delta_{ij}s)^2]$ is the within-cell
dispersion of the score. Dividing by $V$ makes the penalties scale-free, so shrinking the
scores cannot satisfy them; a weight is zero when that kind of pair is unavailable. The
criterion uses the labels only through $y_{ij}$, so it is unchanged by any strictly increasing
relabeling within a cell: it never sees $\psi_{ct}$.

Only the first term depends on the error assumption. Under median zero an ordinal loss makes
sense: the probability of misordering a pair,
$Q^{\mathrm{cross}}_{\mathrm a}(s)=\Pr\bigl((2y_{ij}-1)\Delta_{ij}s<0\bigr)+\tfrac12\Pr(\Delta_{ij}s=0)$,
trained through a smooth surrogate. Under heterogeneous logistic errors it is a RankNet loss
parametrized at the city-year level,
$Q^{\mathrm{cross}}_{\mathrm b}(s,\eta)=\E[\ell(\rho_{ct}\Delta_{ij}s,\,y_{ij})]$ with
$\ell(z,y)=\log(1+e^{z})-yz$, $\eta=\{\rho_{ct}\}$ and the scale of $s$ fixed by $V(s)=1$. Under
homogeneous logistic errors it is the RankNet loss of \citet{burges2005},
$Q^{\mathrm{cross}}_{\mathrm c}(s)=\E[\ell(\Delta_{ij}s,\,y_{ij})]$, with the scale absorbed
into $s$.

% =====================================================================
\section{Related work}
% =====================================================================

Rank estimation removes an unknown monotone transformation by comparing pairs
\citep{han1987}, and \citet{abrevaya2000} lets it differ across groups; both keep a
finite-dimensional index. \citet{wang2024} place a deep network inside a rank objective, with
one transformation and independent observations. In learning to rank, \citet{burges2005}
introduce the pairwise logistic loss, and \citet{duchi2010} show that commonly used surrogates
can be inconsistent for the ranking risk, which is why we compare an ordinal loss with the
logistic ones. The stability terms relate to manifold regularization \citep{belkin2006}, slow
feature analysis \citep{wiskott2002} and self-supervised invariance objectives
\citep{bardes2022,jean2019}; neither a pair penalty nor its normalization is new in isolation.
Invariant risk minimization \citep{arjovsky2019} also seeks one predictor across environments
but takes labels to be on a common scale. Wealth prediction from imagery fits one model per
cell \citep{jean2016} or pools dollars \citep{khachiyan2022}, as in Proposition~\ref{p:mix}.

% =====================================================================
\section{Intended results and evaluation}
% =====================================================================

\paragraph{Theory}
The estimator is $\hat f=s_{\hat\theta}$, where $\hat\theta$ minimizes the sample analogue of
\eqref{eq:crit} over a neural network class. Groups are independent clusters with arbitrary
dependence inside each, and asymptotics are in the number of groups. We intend to prove that
$\hat f$ is consistent for $f$ under each version of \ref{a:err}: first what the minimizers of
the population criterion recover, then uniform convergence of the sample criterion. What
consistency means differs by version. Under (c) the index is recovered up to scale and a
constant per cell (a preliminary proof exists). Under (b) the same holds jointly with the scales
$\rho_{ct}$, whose number grows with the sample: the main technical difficulty. Under (a) only
the ordering is recovered, so excess misranking risk vanishes up to the smoothing error of the
surrogate.

\paragraph{Monte Carlo}
The design follows \citet{morris2019} and \citet{kiviet2012}: aims, data-generating processes,
estimands, methods and performance measures are fixed before any run, every method is applied to
the same simulated datasets, and one factor moves at a time from a common baseline. In that
baseline the covariates have three blocks: signal that enters $f$, acquisition noise that is
redrawn at every date, and a group cue that reveals the group; neither of the last two enters
$f$. Labels are generated at the unit level and then transformed by $\psi_{ct}$, never as
independent pair outcomes. Every $f$ is rescaled so that $V(f)=1$ and the best attainable
misranking rate is equal across designs, which holds the difficulty of the task fixed.

\begin{table}[h]
\centering\footnotesize
\renewcommand{\arraystretch}{1.12}
\begin{tabularx}{\textwidth}{@{}>{\raggedright\arraybackslash}p{0.30\textwidth} Y >{\raggedright\arraybackslash}p{0.22\textwidth}@{}}
\toprule
Question & What varies & What is held fixed \\
\midrule
\textbf{Q1. Efficiency against misspecification of $\psi_{ct}$.} When a parametric correction
$\psi_{ct}^{-1}(W)$ is feasible, what does it gain if correct and lose if wrong?
& Shape of $\psi_{ct}$ along one dial (common, cell multiplier, location and scale, shape);
cross-cell differences in the distribution of $f(X)$; number of labeled groups. Competitors:
pooled levels, logs with cell effects, within-cell $z$-scores and percentile ranks, and an
oracle fitted to the latent index.
& $f$; links; errors at (c), where our loss and a correct level regression are both well
specified. \\
\addlinespace
\textbf{Q2. Misspecified errors.} How does each of the three losses behave when its version of
\ref{a:err} fails?
& One departure at a time, each with one dial: tail weight of $e_i-e_j$; dispersion of
$\rho_{ct}$ across cells; dependence of the noise scale on $X$; skewness with a scale that
depends on $X$, which violates (a).
& $\psi_{ct}$; $f$; sample size; links. \\
\addlinespace
\textbf{Q3. Value of stable pairs.} Within each loss, do the links improve on a well-trained
fit without them?
& Link arm: none, temporal only, cross-group only, both. Strength of the acquisition noise and
of the group cue.
& Labels; architecture; tuning budget; errors; $\psi_{ct}$. \\
\addlinespace
\textbf{Q4. Approximate and false links.} How fast do gains erode, and does tuning the
penalty weight prevent harm?
& A drift in $f$ on every link; a share of entirely false links, random or one-sided; penalty
weight fixed or chosen on held-out ranking loss.
& The design of Q3 with both kinds of link. \\
\bottomrule
\end{tabularx}
\caption{Monte Carlo questions, each run for the three losses and for a linear and a nonlinear
$f$.}
\label{tab:mc}
\end{table}

Q1 and Q3 are read from held-out error against the number of labeled groups: a correct
competitor appears as a horizontal gap, the extra labels our method needs, and a misspecified
one as a floor that more data does not remove. Q2 and Q4 are read from error against the dial.

\paragraph{Performance measures}
All are computed in held-out groups and reported with Monte Carlo standard errors: (i) the share
of within-cell pairs ordered differently from the true $f$, comparable across the three losses;
(ii) the same share for pairs across dates and across groups, which a simulation allows because
$f$ is known; and (iii) error in within-cell index differences with $V=1$, for the logistic
losses only, since the ordinal loss identifies no spacing.

% =====================================================================
\section{Execution plan and expected contribution}
% =====================================================================

The work has five steps. (1) Derive the population result for the three versions of
\ref{a:err}. (2) Build the simulator and the three losses from scratch, including the smooth
ordinal surrogate and the minibatch estimate of the ratio penalties, with a linear head on fixed
features and with a small network. (3) Run Q3 first: if links do not help against a well-trained
fit without them, that is the finding we report. (4) Run Q1, Q2 and Q4. (5) Prove sample
consistency for the versions that survive and formalize the regularities the simulations show.
Time permitting, we illustrate the method on U.S.\ aerial imagery and tract incomes.

We expect three contributions: a training criterion that learns one predictor from labels
comparable only within groups, invariant to the label scale by construction; consistency results
that state what each error assumption identifies, from an ordering to an index with
cell-specific reliabilities; and evidence on when unlabeled stable pairs improve prediction,
when they harm it, and what is lost relative to a correct parametric correction.

% =====================================================================
\begin{thebibliography}{99}\small

\bibitem[Abrevaya(2000)]{abrevaya2000}
Abrevaya, J. (2000). Rank estimation of a generalized fixed-effects regression model.
\emph{Journal of Econometrics}, 95(1):1--23.

\bibitem[Arjovsky et~al.(2019)]{arjovsky2019}
Arjovsky, M., Bottou, L., Gulrajani, I., and Lopez-Paz, D. (2019). Invariant risk minimization.
Working paper 1907.02893, arXiv.

\bibitem[Bardes et~al.(2022)]{bardes2022}
Bardes, A., Ponce, J., and LeCun, Y. (2022). VICReg: Variance-invariance-covariance
regularization for self-supervised learning. In \emph{International Conference on Learning
Representations (ICLR)}.

\bibitem[Belkin et~al.(2006)]{belkin2006}
Belkin, M., Niyogi, P., and Sindhwani, V. (2006). Manifold regularization: A geometric framework
for learning from labeled and unlabeled examples. \emph{Journal of Machine Learning Research},
7:2399--2434.

\bibitem[Burges et~al.(2005)]{burges2005}
Burges, C. J.~C., Shaked, T., Renshaw, E., Lazier, A., Deeds, M., Hamilton, N., and Hullender, G.
(2005). Learning to rank using gradient descent. In \emph{Proceedings of the 22nd International
Conference on Machine Learning (ICML)}, pages 89--96.

\bibitem[Duchi et~al.(2010)]{duchi2010}
Duchi, J.~C., Mackey, L.~W., and Jordan, M.~I. (2010). On the consistency of ranking algorithms.
In \emph{Proceedings of the 27th International Conference on Machine Learning (ICML)}.

\bibitem[Han(1987)]{han1987}
Han, A.~K. (1987). Non-parametric analysis of a generalized regression model: The maximum rank
correlation estimator. \emph{Journal of Econometrics}, 35(2--3):303--316.

\bibitem[Jean et~al.(2016)]{jean2016}
Jean, N., Burke, M., Xie, M., Davis, W.~M., Lobell, D.~B., and Ermon, S. (2016). Combining
satellite imagery and machine learning to predict poverty. \emph{Science}, 353(6301):790--794.

\bibitem[Jean et~al.(2019)]{jean2019}
Jean, N., Wang, S., Samar, A., Azzari, G., Lobell, D., and Ermon, S. (2019). Tile2Vec:
Unsupervised representation learning for spatially distributed data. In \emph{Proceedings of the
AAAI Conference on Artificial Intelligence}, 33.

\bibitem[Khachiyan et~al.(2022)]{khachiyan2022}
Khachiyan, A., Thomas, A., Zhou, H., Hanson, G., Cloninger, A., Rosing, T., and Khandelwal,
A.~K. (2022). Using neural networks to predict microspatial economic growth. \emph{American
Economic Review: Insights}, 4(4):491--506.

\bibitem[Kiviet(2012)]{kiviet2012}
Kiviet, J.~F. (2012). Monte Carlo simulation for econometricians. \emph{Foundations and Trends in
Econometrics}, 5(1--2):1--181.

\bibitem[Morris et~al.(2019)]{morris2019}
Morris, T.~P., White, I.~R., and Crowther, M.~J. (2019). Using simulation studies to evaluate
statistical methods. \emph{Statistics in Medicine}, 38(11):2074--2102.

\bibitem[Wang et~al.(2024)]{wang2024}
Wang, T., Zhang, S., Zhang, S., Huang, J., and Ma, S. (2024). Deep transformation model.
Working paper 2410.19226, arXiv.

\bibitem[Wiskott and Sejnowski(2002)]{wiskott2002}
Wiskott, L. and Sejnowski, T.~J. (2002). Slow feature analysis: Unsupervised learning of
invariances. \emph{Neural Computation}, 14(4):715--770.

\end{thebibliography}

\end{document}
